\documentclass[12pt]{article}
\usepackage[UTF8]{ctex}
\usepackage{amsmath,amssymb}
\usepackage{braket}
\usepackage{fancyhdr}
\usepackage{geometry}

% 页面设置
\geometry{a4paper, margin=1in}

% 页眉页脚设置
\pagestyle{fancy}
\fancyhf{}
\fancyhead[L]{量子计算与机器学习第四章作业}
\fancyhead[R]{\today}
\fancyfoot[C]{\thepage}

% 姓名与日期
\title{\textbf{量子计算与机器学习第四章作业}}
\author{姓名：郝思粤\quad 学号：PB23151779}
\date{\today}

\begin{document}

\maketitle

\textbf{题目 4.1}  
证明贝尔态
\[
\ket{\Phi^+}=\frac{1}{\sqrt{2}}\bigl(\ket{00}+\ket{11}\bigr)
\]
可以等效写成
\[
\ket{\Phi^+}=\frac{1}{\sqrt{2}}\bigl(\ket{aa}+\ket{bb}\bigr),
\]
其中 $\ket{a}$ 和 $\ket{b}$ 是任意一组单比特正交归一基。

\bigskip

\textbf{解：}

设单比特计算基为 $\{\ket{0},\ket{1}\}$。  
给定任意一组正交归一基 $\{\ket{a},\ket{b}\}$，根据线性代数知识，必存在一个酉矩阵 $U$ 满足
\[
U\ket{0}=\ket{a},\qquad U\ket{1}=\ket{b}.
\]

考虑两比特系统上的酉算符 $U\otimes U$，则
\begin{align}
(U\otimes U)\ket{\Phi^+}
&=(U\otimes U)\frac{1}{\sqrt{2}}\bigl(\ket{00}+\ket{11}\bigr) \\
&=\frac{1}{\sqrt{2}}\bigl(U\ket{0}\otimes U\ket{0}
  +U\ket{1}\otimes U\ket{1}\bigr) \\
&=\frac{1}{\sqrt{2}}\bigl(\ket{a}\otimes\ket{a}
  +\ket{b}\otimes\ket{b}\bigr) \\
&=\frac{1}{\sqrt{2}}\bigl(\ket{aa}+\ket{bb}\bigr).
\end{align}

因此有
\[
\frac{1}{\sqrt{2}}\bigl(\ket{aa}+\ket{bb}\bigr)
= (U\otimes U)\ket{\Phi^+}.
\]

即贝尔态在任意单比特正交归一基下均可等效写成
\[
\ket{\Phi^+}=\frac{1}{\sqrt{2}}\bigl(\ket{aa}+\ket{bb}\bigr),
\]
证毕。

\end{document}
